% Paso 3: Hoja de Análisis Gramatical
% Proyecto: Análisis Exegético Profundo
% Ministerio "La Gloria es del Señor"

\documentclass[11pt,a4paper]{article}
\usepackage[utf8]{inputenc}
\usepackage[spanish]{babel}
\usepackage[margin=2.5cm]{geometry}
\usepackage{graphicx}
\usepackage{fancyhdr}
\usepackage{xcolor}
\usepackage{tcolorbox}
\usepackage{enumitem}
\usepackage{tabularx}
\usepackage{multirow}
\usepackage{amsfonts}
\usepackage{fontspec}
\usepackage{titlesec}
\usepackage{forest}
\usepackage{tikz}

% Configuración de colores
\definecolor{forestgreen}{RGB}{27,67,50}
\definecolor{sagegreen}{RGB}{45,90,39}
\definecolor{olivegreen}{RGB}{64,145,108}
\definecolor{mintgreen}{RGB}{82,183,136}
\definecolor{lightmint}{RGB}{116,198,157}
\definecolor{softcream}{RGB}{216,243,220}
\definecolor{papergray}{RGB}{118,128,150}

% Configuración de encabezados
\pagestyle{fancy}
\fancyhf{}
\fancyhead[L]{\textcolor{forestgreen}{\textbf{Análisis Exegético Profundo}}}
\fancyhead[R]{\textcolor{papergray}{Paso 3: Análisis Gramatical}}
\fancyfoot[C]{\textcolor{papergray}{\thepage}}
\fancyfoot[R]{\textcolor{papergray}{\scriptsize Ministerio "La Gloria es del Señor"}}

% Configuración de títulos
\titleformat{\section}
{\color{forestgreen}\Large\bfseries}
{\thesection}{1em}{}

\titleformat{\subsection}
{\color{sagegreen}\large\bfseries}
{\thesubsection}{1em}{}

% Configuración de cajas
\tcbuselibrary{skins}
\newtcolorbox{infobox}[1]{
    colback=softcream,
    colframe=olivegreen,
    fonttitle=\bfseries,
    title=#1,
    left=10pt,
    right=10pt,
    top=10pt,
    bottom=10pt
}

\newtcolorbox{checklistbox}{
    colback=white,
    colframe=mintgreen,
    left=15pt,
    right=15pt,
    top=10pt,
    bottom=10pt,
    boxrule=2pt
}

\newtcolorbox{grammarbox}{
    colback=lightmint!20,
    colframe=sagegreen,
    left=10pt,
    right=10pt,
    top=8pt,
    bottom=8pt,
    boxrule=1.5pt
}

\begin{document}

% Título principal
\begin{center}
    {\Huge\color{forestgreen}\textbf{PASO 3}}\\[0.5cm]
    {\Large\color{sagegreen}\textbf{ANÁLISIS GRAMATICAL}}\\[0.3cm]
    {\large\color{papergray}Estructura Sintáctica y Relaciones Gramaticales}\\[1cm]
    
    \begin{tcolorbox}[colback=olivegreen!10, colframe=olivegreen, width=0.8\textwidth]
        \centering
        \textcolor{forestgreen}{\textbf{Proyecto: Análisis Exegético Profundo}}\\
        \textcolor{papergray}{Metodología de 8 Pasos para Estudio Bíblico Riguroso}
    \end{tcolorbox}
\end{center}

\vspace{1cm}

% Información del estudiante
\begin{infobox}{Información del Proyecto}
\begin{tabularx}{\textwidth}{|X|X|}
\hline
\textbf{Nombre del Estudiante:} & \rule{6cm}{0.4pt} \\[0.8cm]
\hline
\textbf{Texto Bíblico:} & \rule{6cm}{0.4pt} \\[0.8cm]
\hline
\textbf{Idioma Original:} & $\square$ Hebreo \quad $\square$ Griego \quad $\square$ Arameo \\[0.8cm]
\hline
\textbf{Fecha:} & \rule{6cm}{0.4pt} \\[0.8cm]
\hline
\end{tabularx}
\end{infobox}

\section{Objetivo del Paso 3: Análisis Gramatical}

El análisis gramatical examina la estructura sintáctica del texto para entender cómo las palabras se relacionan entre sí y cómo contribuyen al mensaje del autor. Este análisis revela el flujo lógico del pensamiento y las relaciones de significado.

\subsection{Descripción del Proceso}

Examinaremos la estructura gramatical y sintáctica del texto, identificando sujetos, verbos, objetos y modificadores. Analizaremos tiempo, modo y voz de los verbos, así como las relaciones entre cláusulas y el flujo argumentativo.

\begin{infobox}{Herramientas de Análisis}
\textbf{Para Griego:} Análisis morfológico, diagramas de oración, sintaxis del griego koiné\\
\textbf{Para Hebreo:} Análisis verbal, sintaxis hebrea, patrones poéticos\\
\textbf{Recursos:} Wallace, Dana \& Mantey, Waltke \& O'Connor, BDF
\end{infobox}

\section{Lista de Verificación}

\begin{checklistbox}
\textbf{Marca cada elemento completado:}

\begin{itemize}[leftmargin=1cm]
    \item[$\square$] Analizar formas verbales (tiempo, modo, voz)
    \item[$\square$] Identificar casos de sustantivos y pronombres
    \item[$\square$] Examinar relaciones entre cláusulas
    \item[$\square$] Determinar el flujo lógico del argumento
    \item[$\square$] Crear diagrama de oración (opcional pero recomendado)
    \item[$\square$] Identificar construcciones sintácticas especiales
    \item[$\square$] Analizar el uso de partículas y conectores
    \item[$\square$] Evaluar la estructura retórica del pasaje
\end{itemize}
\end{checklistbox}

\newpage

\section{Análisis Verbal}

\subsection{Inventario de Verbos Principales}

\begin{tabularx}{\textwidth}{|c|X|c|c|c|c|X|}
\hline
\textbf{Vers.} & \textbf{Verbo} & \textbf{Tiempo} & \textbf{Modo} & \textbf{Voz} & \textbf{Persona} & \textbf{Significado Gramatical} \\
\hline
& & & & & & \\[1cm]
\hline
& & & & & & \\[1cm]
\hline
& & & & & & \\[1cm]
\hline
& & & & & & \\[1cm]
\hline
& & & & & & \\[1cm]
\hline
& & & & & & \\[1cm]
\hline
& & & & & & \\[1cm]
\hline
& & & & & & \\[1cm]
\hline
\end{tabularx}

\subsection{Análisis del Sistema Verbal}

\textbf{¿Qué tiempos verbales predominan y qué indica esto sobre la naturaleza del texto?}
\vspace{2cm}

\textbf{¿Hay algún cambio significativo en tiempos/modos que indique estructura o énfasis?}
\vspace{2cm}

\textbf{¿Cómo contribuyen los aspectos verbales al mensaje del pasaje?}
\vspace{2cm}

\section{Análisis de Casos y Sintaxis}

\subsection{Análisis de Sustantivos y Casos}

\begin{grammarbox}
\textbf{Griego:} Nominativo, Genitivo, Dativo, Acusativo, Vocativo\\
\textbf{Hebreo:} Estado constructo, absoluto, función sintáctica
\end{grammarbox}

\begin{tabularx}{\textwidth}{|c|X|c|X|X|}
\hline
\textbf{Vers.} & \textbf{Sustantivo/Pronombre} & \textbf{Caso} & \textbf{Función} & \textbf{Significado Sintáctico} \\
\hline
& & & & \\[1cm]
\hline
& & & & \\[1cm]
\hline
& & & & \\[1cm]
\hline
& & & & \\[1cm]
\hline
& & & & \\[1cm]
\hline
& & & & \\[1cm]
\hline
\end{tabularx}

\newpage

\subsection{Relaciones entre Cláusulas}

\textbf{Identifica las cláusulas principales y subordinadas:}

\begin{tabularx}{\textwidth}{|c|X|X|X|}
\hline
\textbf{Vers.} & \textbf{Cláusula} & \textbf{Tipo} & \textbf{Relación con el contexto} \\
\hline
& & Principal / Subordinada & \\[1cm]
\hline
& & Principal / Subordinada & \\[1cm]
\hline
& & Principal / Subordinada & \\[1cm]
\hline
& & Principal / Subordinada & \\[1cm]
\hline
& & Principal / Subordinada & \\[1cm]
\hline
\end{tabularx}

\subsection{Conectores y Partículas}

\begin{tabularx}{\textwidth}{|c|X|X|X|}
\hline
\textbf{Vers.} & \textbf{Conector/Partícula} & \textbf{Función Gramatical} & \textbf{Efecto en el Argumento} \\
\hline
& & & \\[1cm]
\hline
& & & \\[1cm]
\hline
& & & \\[1cm]
\hline
& & & \\[1cm]
\hline
\end{tabularx}

\section{Estructura del Argumento}

\subsection{Flujo Lógico}

\textbf{¿Cuál es la progresión lógica del argumento en el pasaje?}
\vspace{3cm}

\textbf{¿Qué palabras o construcciones indican transiciones importantes?}
\vspace{2cm}

\textbf{¿Hay estructura quiástica, paralela u otros patrones literarios?}
\vspace{2cm}

\subsection{Construcciones Sintácticas Especiales}

\textbf{Genitivos especiales (si aplica):}
\vspace{1.5cm}

\textbf{Construcciones infinitivas:}
\vspace{1.5cm}

\textbf{Construcciones participiales:}
\vspace{1.5cm}

\textbf{Discurso directo/indirecto:}
\vspace{1.5cm}

\newpage

\section{Diagrama de Oración (Opcional)}

\begin{infobox}{Instrucciones para el Diagrama}
Crea un diagrama de oración del versículo más importante o complejo del pasaje. Esto te ayudará a visualizar las relaciones gramaticales.
\end{infobox}

\textbf{Versículo seleccionado para diagramar:} \rule{8cm}{0.4pt}

\vspace{6cm}
\begin{center}
[Espacio para el diagrama de oración]
\end{center}
\vspace{4cm}

\textbf{Explicación del diagrama:}
\vspace{3cm}

\section{Análisis de Estructura Retórica}

\subsection{Organización del Material}

\textbf{¿Cómo está organizado el material en el pasaje?}
\begin{itemize}[leftmargin=1cm]
    \item[$\square$] Cronológicamente
    \item[$\square$] Lógicamente (causa-efecto)
    \item[$\square$] Por comparación/contraste
    \item[$\square$] Por importancia
    \item[$\square$] Poéticamente
    \item[$\square$] Otro: \rule{6cm}{0.4pt}
\end{itemize}

\subsection{Énfasis Gramatical}

\textbf{¿Qué elementos reciben énfasis especial por su posición o construcción?}
\vspace{3cm}

\textbf{¿Hay repeticiones o variaciones gramaticales significativas?}
\vspace{2cm}

\section{Síntesis Gramatical}

\subsection{Hallazgos Principales}

\textbf{¿Qué revelan las estructuras gramaticales sobre el mensaje del autor?}
\vspace{3cm}

\textbf{¿Hay algún aspecto gramatical que afecte significativamente la interpretación?}
\vspace{3cm}

\textbf{¿Cómo contribuye el análisis gramatical a tu comprensión del flujo del argumento?}
\vspace{3cm}

\begin{infobox}{Principio Hermenéutico}
La gramática es la sirvienta de la exégesis, no su señora. Busca patrones significativos pero no fuerces interpretaciones basadas únicamente en la gramática.
\end{infobox}

\section{Siguiente Paso}

Con el análisis gramatical completado, procederás al \textbf{Paso 4: Contexto Histórico-Cultural}, donde investigaremos el trasfondo del mundo bíblico.

\vspace{1cm}

\begin{center}
\begin{tcolorbox}[colback=mintgreen!20, colframe=mintgreen, width=0.7\textwidth]
\centering
\textcolor{forestgreen}{\textbf{¡Análisis Preciso!}}\\
\textcolor{papergray}{Has examinado la estructura gramatical del texto con rigor académico.}
\end{tcolorbox}
\end{center}

\end{document}